\documentclass[Chp3.tex]{subfiles}

\begin{document}


\section{Proposing a measure of wealth inequality aversion}
\label{sec:Altmeasure}


\par The ranking over wealth distributions via specifications on $U(y)$ will allow us to propose measures of inequality. To see the novelty in this \say{social welfare approach}, notice that measures of inequality are generally statistical objects. Namely, the variance, coefficient of variation, and mean deviation are each calculated using collected data on the distribution of an outcome such as income or wealth.

\par Consider one plausible inequality measure: given some present distribution $f(w)=\{p_s(y)\}_{s\in S}$, the ratio between social welfare at the current distribution and social welfare at equal wealth $\mu$. Using the robust welfare index
\[
W^{A} \equiv \min_{\lambda \in C}\int_S\int_0^{\bar y} U(y)\,p_s(y)\,dy\,d\lambda,
\qquad C\subseteq \Delta(S),
\]
we may write
\[
D^{A} = \frac{W^{A}}{U(\mu)}.
\]

\par \textit{Notation.} Superscript $A$ denotes ambiguity objects; this is the first appearance of $D^{A}$ in the main text.

\par However, this ratio is not invariant to positive affine transformations of $U$, so different cardinal representations can produce different values of $D^{A}$ even when they induce the same ranking.\footnote{As shown by \cite{hd20} and others.}

\subsection{The equally distributed equivalent measure of inequality}

\par To avoid this undesirable property, define the equally distributed equivalent wealth level as follows.

\begin{dfn}
The equally distributed equivalent level of wealth $w_{EDE}$ is defined by
\[
U(w_{EDE}) = W^{A}.
\]
Equivalently,
\[
U(w_{EDE}) = \min_{\lambda \in C}\int_S\int_0^{\bar y} U(y)\,p_s(y)\,dy\,d\lambda.
\]
\end{dfn}

\par Since each $p_s$ integrates to one on $Y$ and each $\lambda\in C$ is a probability on $S$, the normalization factor is one, so no extra scale term is needed in the definition.

\par The corresponding equally distributed equivalent inequality measure is
\[
I^{A} = 1 - \frac{w_{EDE}}{\mu}.
\]

\par \textit{Notation.} This is the first appearance of $I^{A}$ in the main text.

\par This measure is cardinally robust (with respect to positive affine transformations of $U$) and lies in $[0,1]$ under the usual ordering conditions.

\subsection{Closed-form expression under homothetic preferences}

\par For computation, the key identity is
\[
w_{EDE} = U^{-1}(W^{A}).
\]
For the homothetic class,
\[
U(y)=
\begin{cases}
\alpha + \beta\dfrac{y^{1-\epsilon}}{1-\epsilon}, & \epsilon\neq 1,\\
\alpha + \beta\ln y, & \epsilon = 1,
\end{cases}
\qquad \beta>0,
\]
with $y>0$ and $\epsilon>0$ for strict concavity (note: $\epsilon=0$ is linear and therefore not strictly concave).

\par If $\epsilon\neq 1$,
\[
w_{EDE} = \left[\frac{1-\epsilon}{\beta}\,(W^{A}-\alpha)\right]^{\frac{1}{1-\epsilon}},
\qquad
I^{A} = 1 - \frac{1}{\mu}\left[\frac{1-\epsilon}{\beta}\,(W^{A}-\alpha)\right]^{\frac{1}{1-\epsilon}}.
\]

\par If $\epsilon=1$,
\[
w_{EDE} = \exp\!\left(\frac{W^{A}-\alpha}{\beta}\right),
\qquad
I^{A} = 1 - \frac{1}{\mu}\exp\!\left(\frac{W^{A}-\alpha}{\beta}\right).
\]

\par Appendix~\ref{app:objective-benchmark} gives the same derivation for the objective-uncertainty benchmark and then shows that the ambiguity case is obtained by replacing $W^{O}$ with $W^{A}$.

\subsection{Wealth inequality aversion and the racial wealth gap}
\label{subsec:compute-measures}

\subsubsection{Data}

\par I use the triennial Survey of Consumer Finances (SCF), $1989$--$2022$, and restrict the sample to Black and White households in each wave. For each year, I compute weighted medians for net wealth and gross assets, and then compute $I^{O}$ and $I^{A}$ under two treatments: (i) a restricted sample that keeps only strictly positive wealth values, and (ii) a shifted full sample that adds a year-specific constant so all observations are strictly positive.

\par The ambiguity aggregator is computed under $C=\Delta(S)$ with $S=\{\text{Black},\text{White}\}$, so the maxmin benchmark uses the worst-performing racial state each year. The main text reports $\epsilon\in\{0.25,0.5,1,1.5,2\}$; the linear benchmark $\epsilon=0$ is discussed in Appendix~\ref{app:eps0-benchmark}.

\par With these modeling objects fixed, the theoretical welfare expressions map directly into computable objects in the data: for each year and wealth concept, we compute $W^{O}$ and $W^{A}$, invert $U$ to recover the equally distributed equivalent wealth level, and then construct $I^{O}$ and $I^{A}$. The remaining step is to set values for $\epsilon$, which governs the curvature of $U(y)$ and therefore the degree of inequality aversion imposed in the computation.

\par The parameter $\epsilon$ indexes the curvature of the social welfare function $U(y)$ and therefore the strength of inequality aversion: larger $\epsilon$ means stronger concavity, a higher penalty on low-wealth outcomes, and thus higher implied inequality aversion in both $I^{O}$ and $I^{A}$.

\subsubsection{The Black-White gap in net wealth and gross assets}

\begin{figure}[H]
\centering
\includegraphics[width=0.92\linewidth]{\subfix{../Figures/03_gap_diagnostic.png}}
\caption{Black-White median wealth differential in net wealth and gross assets, SCF 1989--2022.}
\label{fig:gap-net-gross}
\end{figure}

\par Figure~\ref{fig:gap-net-gross} shows that the median Black-White differential is positive and persistent in both series, with gross-asset gaps generally larger than net-wealth gaps over the sample.

\par Checking gross assets alongside net wealth is useful because it verifies that the empirical trend is not an artifact of liability measurement alone: both series move with similar broad dynamics over time. At the same time, the gross-asset gap is higher in levels than the net-wealth gap, which is consistent with definition since net wealth subtracts debts while gross assets do not.


\subsubsection{$\epsilon=0.25$}

\begin{figure}[H]
\centering
\includegraphics[width=0.92\linewidth]{\subfix{../Figures/08_measures_eps025_net.png}}\\[0.75em]
\includegraphics[width=0.92\linewidth]{\subfix{../Figures/08_measures_eps025_gross.png}}
\caption{Inequality measures for $\epsilon=0.25$ (top: net wealth, bottom: gross assets). In each panel pair, shifted full sample is on the left and restricted sample is on the right.}
\label{fig:eps025-measures}
\end{figure}

\par The case $\epsilon=0.25$ is closest to the linear benchmark among the reported nonlinear specifications, so curvature in $U(y)$ is relatively mild.

\par Even with this low curvature, the ambiguity ranking already yields higher inequality levels than the objective ranking in both wealth concepts (\(I^{A}>I^{O}\)); the restricted sample is higher and smoother over time, while the shifted full sample is more cyclical.

\par Relative to the higher-curvature specifications ($\epsilon=0.5,1,1.5,2$), the $\epsilon=0.25$ series are lower in level, which is exactly what we should expect when inequality aversion is weaker.

\subsubsection{$\epsilon=0.5$}

\begin{figure}[H]
\centering
\includegraphics[width=0.92\linewidth]{\subfix{../Figures/05_measures_eps05_net.png}}\\[0.75em]
\includegraphics[width=0.92\linewidth]{\subfix{../Figures/05_measures_eps05_gross.png}}
\caption{Inequality measures for $\epsilon=0.5$ (top: net wealth, bottom: gross assets). In each panel pair, shifted full sample is on the left and restricted sample is on the right.}
\label{fig:eps05-measures}
\end{figure}

\par The case $\epsilon=0.5$ corresponds to intermediate curvature in $U(y)$: stronger inequality aversion than $\epsilon=0.25$, but weaker than the log case $\epsilon=1$.

\par For $\epsilon=0.5$, the restricted sample produces high and comparatively stable inequality levels over time, while the shifted full sample varies more across years. In both net wealth and gross assets, $I^{A}$ is consistently above $I^{O}$, so the objective benchmark understates measured inequality relative to the ambiguity benchmark.

\subsubsection{$\epsilon=1$}

\begin{figure}[H]
\centering
\includegraphics[width=0.92\linewidth]{\subfix{../Figures/06_measures_eps1_net.png}}\\[0.75em]
\includegraphics[width=0.92\linewidth]{\subfix{../Figures/06_measures_eps1_gross.png}}
\caption{Inequality measures for $\epsilon=1$ (top: net wealth, bottom: gross assets). In each panel pair, shifted full sample is on the left and restricted sample is on the right.}
\label{fig:eps1-measures}
\end{figure}

\par The case $\epsilon=1$ is the log-utility benchmark, which increases curvature relative to $\epsilon=0.5$ and therefore strengthens inequality aversion.

\par The same ranking persists at $\epsilon=1$: in every case, $I^{A}>I^{O}$. Relative to $\epsilon=0.5$, the restricted-sample indices move upward in both wealth concepts, while remaining smoother over time than the shifted full-sample series.


\subsubsection{$\epsilon=1.5$}

\begin{figure}[H]
\centering
\includegraphics[width=0.92\linewidth]{\subfix{../Figures/09_measures_eps15_net.png}}\\[0.75em]
\includegraphics[width=0.92\linewidth]{\subfix{../Figures/09_measures_eps15_gross.png}}
\caption{Inequality measures for $\epsilon=1.5$ (top: net wealth, bottom: gross assets). In each panel pair, shifted full sample is on the left and restricted sample is on the right.}
\label{fig:eps15-measures}
\end{figure}

\par The case $\epsilon=1.5$ adds substantial curvature beyond the log benchmark, so inequality aversion is stronger than at $\epsilon=1$ but weaker than at $\epsilon=2$.

\par In the shifted full sample, this is where the pull-up becomes clear: levels rise materially relative to $\epsilon\in\{0.25,0.5,1\}$, although the series is still not uniformly near one in all years. In the restricted sample, both indices remain very high and smooth, with $I^{A}>I^{O}$ preserved.

\subsubsection{$\epsilon=2$}

\begin{figure}[H]
\centering
\includegraphics[width=0.92\linewidth]{\subfix{../Figures/07_measures_eps2_net.png}}\\[0.75em]
\includegraphics[width=0.92\linewidth]{\subfix{../Figures/07_measures_eps2_gross.png}}
\caption{Inequality measures for $\epsilon=2$ (top: net wealth, bottom: gross assets). In each panel pair, shifted full sample is on the left and restricted sample is on the right.}
\label{fig:eps2-measures}
\end{figure}

\par The case $\epsilon=2$ imposes strong curvature in $U(y)$ and thus the strongest inequality aversion among the reported specifications.

\par As $\epsilon$ rises from $1.5$ to $2$, restricted-sample inequality indices concentrate even more tightly near one, reflecting the stronger curvature penalty on lower-wealth outcomes. Even in this high-curvature case, the ambiguity measure remains weakly above the objective measure, preserving the core result $I^{A}\ge I^{O}$ throughout.

\par Taken together, the inequality-measure series tell the same within-sample story across specifications (from $\epsilon=0.25$ through $\epsilon=2$): inequality remains persistently high over time, with fluctuations but no sustained convergence. This is consistent with Figure~\ref{fig:gap-net-gross}, which shows a persistently positive Black-White median wealth differential throughout the sample.

\par In the restricted positive-only sample, measured inequality is especially robust to specification in the sense that levels remain high across all reported $\epsilon$, and the indices move close to one as curvature rises (most clearly for $\epsilon\ge 1$, and for both $I^{O}$ and $I^{A}$ at $\epsilon=2$). This pattern is therefore not only a maxmin artifact: maxmin mainly raises levels via the $I^{A}>I^{O}$ wedge, while stronger curvature itself also pushes indices upward. By contrast, in the shifted full-sample treatment both $I^{O}$ and $I^{A}$ are more variable over time and show episodes of comparatively low measured inequality at low and moderate $\epsilon$; they begin to pull upward noticeably at $\epsilon=1.5$, and become frequently near one only in the highest-curvature case $\epsilon=2$. This higher time variation in the shifted series tracks the median-gap dynamics in Figure~\ref{fig:gap-net-gross} more closely than the restricted-sample series. Appendix Figure~\ref{fig:gap-net-gross-restricted} shows that, even under the same positive-only restriction, both net-wealth and gross-assets median gaps still vary over time and trend upward; relative to that diagnostic, the restricted-sample theoretical indices are unusually stable and persistently high, with $I^{A}$ near the top throughout and $I^{O}$ increasing as $\epsilon$ increases.




\end{document}
